% Canonical Paper I source.

The integral pillar has two layers.  In one dimension it is the mixed
additive--multiplicative integral whose two sides are Scott-continuous in the
two directions of Section~\ref{sec:projective-affine}; in higher dimension it is
the order defect as a group-valued $1$-cocycle on the linear-extension graph,
closed by the holographic Stokes theorem.  Throughout, $u,v$ are densities on
$[a,b]$ and $M(t)=\int_a^t v$ is the multiplicative charge.

\subsection{The one-dimensional mixed affine integral}

\begin{definition}[Mixed affine integral]
\label{def:mixed-affine-integral}
For $u,v\in L^1([a,b])$, define the \emph{continuous mixed additive--multiplicative
integral} in absolute-time parametrization,
\[
  \mathfrak M_a^b[u;v]:=\int_a^b u(x)\,\exp\!\Bigl(\int_x^b v(s)\,ds\Bigr)\,dx ,
\]
and, in the right-endpoint sliding-window parametrization, the one-sided
Volterra convolution
\[
  \mathfrak A(w):=\int_0^w u(x)\,K(w-x)\,dx,
  \qquad
  K(r)=\exp\!\Bigl(\int_0^r v(y)\,dy\Bigr).
\]
\end{definition}

\begin{proposition}[Mixed-integral equivalences]
\label{prop:mixed-integral-equivalences}
\begin{enumerate}
  \item \textbf{Affine transport and splitting.} $\mathfrak M_a^b[u;v]$ is the
  inhomogeneous propagation term of the affine transport equation
  $w'(t)=v(t)w(t)+u(t)$:
  \[
    w(b)=e^{\int_a^b v}w(a)+\mathfrak M_a^b[u;v],
  \]
  and satisfies the splitting law and the right-endpoint derivative
  \[
    \mathfrak M_a^c[u;v]=e^{\int_b^c v}\mathfrak M_a^b[u;v]+\mathfrak M_b^c[u;v],
    \qquad
    \frac{d}{db}\mathfrak M_a^b[u;v]=u(b)+v(b)\,\mathfrak M_a^b[u;v].
  \]
  \item \textbf{Window derivative and integration by parts.} For
  $u\in C^1([0,w])$, $v\in C([0,w])$,
  \[
    \mathfrak A'(w)=u(0)K(w)+\int_0^w u'(x)K(w-x)\,dx
    =u(w)+\int_0^w u(x)\,v(w-x)K(w-x)\,dx .
  \]
  \item \textbf{Reverse ACS line integral.} On the accumulative commutative space
  $\mathcal A=\mathbb R^2_{(A,M)}$,
  \[
    \mathfrak M_a^b[u;v]=\int_{C^{\mathrm{rev}}}\eta ,
    \qquad
    \eta=e^M dA ,
  \]
  where $C^{\mathrm{rev}}$ is the ACS path generated by the reversed density.
  \item \textbf{Twisted Rota--Baxter identity.} With
  $(I_M f)(t)=e^{M(t)}\int_{t_0}^t e^{-M(s)}f(s)\,ds=(e^M\circ J\circ e^{-M})f$
  and the twisted product $\star_M$ of
  Proposition~\ref{prop:graded-leibniz-and-gauge},
  \[
    I_M(f)\star_M I_M(g)=I_M\bigl(I_M(f)\star_M g+f\star_M I_M(g)\bigr).
  \]
\end{enumerate}
\end{proposition}

\begin{proof}
(1) Differentiate $e^{-\int_a^t v}w(t)$ to isolate the inhomogeneous term; the
splitting law is the semigroup property of the fundamental solution, and the
endpoint derivative is the transport equation itself.  (2) Differentiate the
convolution under the sign and integrate by parts.  (3) Substitute $A,M$ and
observe that the density $\eta=e^M dA$ integrates along the reversed path to the
displayed form.  (4) is a direct computation using the twisted derivation law
of Proposition~\ref{prop:graded-leibniz-and-gauge}.
\end{proof}

\subsection{Dual adjoint pairing and the Lawson two-sided functional}

The two directions of the mixed integral are not symmetric: the additive source
is the Scott (from-below) side and the multiplicative density is the lower
(dual) side.  Their pairing is the functional form of the Lawson topology.

\begin{theorem}[Lawson dual adjoint pairing]
\label{thm:lawson-dual-adjoint-pairing}
\begin{enumerate}
  \item \textbf{Dual connection.} On the weight-${(-1)}$ dual line, define
  $\nabla^{(-1)}p:=\delta_H p+p\,dM$.  Then
  \[
    \delta_H(pa)=\bigl(\nabla^{(-1)}p\bigr)a+p\,\bigl(\nabla^{(1)}a\bigr),
  \]
  and when $p$ satisfies the homogeneous backward-parallel equation
  $\nabla^{(-1)}p=0$, i.e.\ $-p'(t)=v(t)p(t)$ with $p(b)=1$, one has
  $\delta_H(pa)=p\,dA$, so that
  \[
    \mathfrak M_a^b[u;v]=\int_a^b p(t)\,u(t)\,dt .
  \]
  \item \textbf{Backward adjoint equals the reciprocal image.} With the forward
  scale $K(t)=e^{M(t)}$ and the backward adjoint
  $p(t)=\exp\!\bigl(\int_t^b v(s)\,ds\bigr)$,
  \[
    p(t)=e^{M(b)-M(t)}=K(b)\cdot\neg\bigl(K(t)\bigr),
    \qquad
    \neg(a)=\frac1a .
  \]
  \item \textbf{Bilateral Scott continuity.} For nonnegative densities
  $u,v\ge0$, the functional $u\mapsto\mathfrak M[u;v]$ and
  $v\mapsto\mathfrak M[u;v]$ are each monotone and preserve directed suprema,
  hence are separately Scott-continuous; the pairing $\int_a^b u\,p$ is the
  functional-level realization of the Lawson topology
  $\lambda=\sigma\vee\neg_*\sigma$ of
  Proposition~\ref{prop:generator-domain-roles}.
\end{enumerate}
\end{theorem}

\begin{proof}
(1) is the graded Leibniz law of
Proposition~\ref{prop:graded-leibniz-and-gauge} at weights $-1,1$; solving
$\nabla^{(-1)}p=0$ gives $p(t)=e^{M(b)-M(t)}$.  (2) is the displayed
computation.  (3) monotone convergence gives the directed-sup preservation in
each argument; the forward scale $K$ and the backward adjoint $p=\neg(K)$ carry
the Scott and the lower side respectively.
\end{proof}

\subsection{The order defect as a group-valued $1$-cocycle}

For a general dependency poset $P$ the order defect between two linear
extensions is additive along adjacent swaps, and closed by the Jacobi identity.

\begin{theorem}[Defect group $1$-cocycle]
\label{thm:defect-group-1-cocycle}
Let $P$ be a finite dependency poset and $G(P)$ its linear-extension graph
(vertices are linear extensions $\gamma$, edges are adjacent transpositions of
incomparable elements).
\begin{enumerate}
  \item \textbf{Exact telescoping.} For a path
  $\gamma_0\to\gamma_1\to\cdots\to\gamma_m$ in $G(P)$,
  \[
    \rho(\gamma_m)\rho(\gamma_0)^{-1}
    =\delta(\gamma_{m-1}\to\gamma_m)\cdots\delta(\gamma_0\to\gamma_1),
    \qquad
    \delta(\gamma\to\gamma'):=\rho(\gamma')\rho(\gamma)^{-1},
  \]
  and the scalar defect $\nu_a(\gamma_0)-\nu_a(\gamma_m)$ is the corresponding
  telescoping sum.  Both depend only on the endpoints.
  \item \textbf{Suffix-conjugated commutator.} In a history $\gamma=UxyW$ ($U$
  executed first, $W$ last), swapping adjacent incomparable $x,y$ to
  $\gamma'=UyxW$ gives
  \[
    \delta(UxyW\to UyxW)=\rho(W)\,\bigl[\rho(x),\rho(y)\bigr]\,\rho(W)^{-1},
    \qquad
    [g,h]:=ghg^{-1}h^{-1} .
  \]
  \item \textbf{BCH expansion and Jacobi closure.} With infinitesimal
  generators $e^{hX},e^{kY},e^{lZ}$,
  \[
    [e^{hX},e^{kY}]
    =\exp\Bigl(hk[X,Y]+\tfrac{h^2k}{2}[X,[X,Y]]+\tfrac{hk^2}{2}[Y,[X,Y]]+\cdots\Bigr),
  \]
  and the two three-step paths around the hexagon of a three-element antichain
  $\{X,Y,Z\}$ differ, to third order, by exactly the Jacobi identity
  $[X,[Y,Z]]+[Y,[Z,X]]+[Z,[X,Y]]=0$ (the discrete Bianchi identity
  $d\Omega=0$).
\end{enumerate}
\end{theorem}

\begin{proof}
(1) is the telescoping product; (2) is the direct computation
$\rho(UyxW)\rho(UxyW)^{-1}=\rho(W)\rho(x)\rho(y)\rho(x)^{-1}\rho(y)^{-1}\rho(W)^{-1}$;
(3) is the Campbell--Baker--Hausdorff formula and the vanishing of its
third-order mixed term by the Jacobi identity.
\end{proof}

\subsection{Holographic Stokes theorem}

The two-dimensional order defect is a boundary quantity: the bulk carries an
exact curvature form, and Stokes pushes its integral to the boundary holonomy.

\begin{proposition}[Holographic Stokes, minimal case]
\label{prop:holographic-stokes-minimal}
In the contact realization of the minimal branching DAG $\{A,M\}$, the curvature
two-form
\[
  d\alpha=-\lambda\,da\wedge dv=d(-\lambda a\,dv)
\]
is nonzero but exact.  For the charge rectangle $\Sigma=[0,h]\times[0,k]$, Stokes
gives
\[
  T=\int_\Sigma d\alpha=\oint_{\partial\Sigma}(-\lambda a\,dv),
\]
i.e.\ the order defect is the boundary holonomy of the potential one-form
$-\lambda a\,dv$ around the closed loop
$\partial\Sigma=(A_h,M_k,A_{-h},M_{-k})$.
\end{proposition}

\begin{proof}
$d\alpha=-\lambda\,da\wedge dv$ is exact with primitive $-\lambda a\,dv$; Stokes
identifies the area integral with the boundary line integral, and the boundary
integral is the closed-loop drift of
Theorem~\ref{thm:dag-defect-equals-curvature}.
\end{proof}

\begin{remark}[The ideal surface as the true limit]
\label{rem:scott-preserves-ideal-surface}
On the completion axis, the finite DAGs (compact elements) form a flat
``bulk'' and the infinite histories $\Gamma^\omega$ form an ``ideal surface''.
The Scott topology keeps the infinite histories as genuine limits, so the
curvature exists as a surface holonomy, whereas the Alexandrov topology would
reify them as isolated points and erase the bulk--surface distinction.  The
macroscopic discrete Stokes integral over an arbitrary dependency poset $P$ and
its spectral docking with the down-set lattice of
Theorem~\ref{thm:birkhoff-downset-spinal} are open (see OQ-095.3, OQ-095.4).
This remark is a structural proposal, not a theorem.
\end{remark}
